% !TeX root = ../../../main.tex
\documentclass[../../../main.tex]{subfiles}
\begin{document}

\subsubsection{Background oriented schlieren (BOS)}

\paragraph{Literature on BOS}

Schlieren imaging encompasses a broad family of techniques, implemented in many different ways, for visualizing the internal structure of transparent flows characterized by non-uniform refractive index fields. This work concentrates on the Background-Oriented Schlieren (BOS) approach, the newest variant within this family and one that continues to drive innovation among experimental researchers. This technique was independently introduced in the peer-reviewed literature in the year 2000: Dalziel and colleagues referred to it as ``synthetic Schlieren''~\cite{Dalziel2000WholefieldDensity}, while Raffel and coworkers coined the term ``Background-Oriented Schlieren''~\cite{Raffel2000ApplicabilityBackground}. The latter designation has since become the standard terminology within the fluid mechanics community.

Compared to other Schlieren-based methods, BOS offers several benefits, chief among them the minimal experimental requirements: only a camera to photograph a textured background and the background pattern itself are needed. A BOS measurement quantifies how the recorded pattern appears distorted from the camera's viewpoint, a distortion caused by light-ray bending due to refractive-index variations occurring in the region between the camera and the background. While a synthetic random pattern is typically projected or printed for use as the background, naturally occurring textures can serve the same purpose such as landscapes~\cite{Hargather2010NaturalbackgroundorientedSchlieren}, and even the granular intensity variations of band-pass-filtered sunlight at the solar surface~\cite{Jr2017BackgroundOriented}. These natural-background applications have notably enabled the visualization of aerodynamic flow fields surrounding helicopters and fighter jets in flight~\cite{AirborneApplication, Heineck2016BackgroundOriented, Schmidt2025TwentyFiveYears}.

Several studies have systematically compared the influence of background-pattern properties (including dot shape, color, and grayscale content) on BOS measurement quality~\cite{Hu2025PerformanceBackgroundoriented, Leopold2012INCREASEACCURAY}. Hu et al.~\cite{Hu2025PerformanceBackgroundoriented} evaluated semi-random dot patterns alongside fractal-like backgrounds generated by diffusion-limited aggregation (dot-DLA), reporting that the conventional dot pattern was the most robust choice for well-defined explosion flows, whereas the dot–DLA pattern performed best at capturing fine-scale turbulent structures. Given its robustness across flow regimes and its comparative simplicity to generate and implement, the dot pattern remains the most widely adopted background in BOS experiments, and it is the pattern retained in the present work.